\section{Weekly Summary}

\begin{tcolorbox}[colback=yellow!10,colframe=black!50,title=AI-Generated Content]
\noindent The summary below was written by Claude (Anthropic) from that week's regression outputs and series descriptions. It has not been reviewed or fact-checked by a human, and should be read as an automated, informal recap -- not analysis.
\end{tcolorbox}

Welcome back to the weekly ritual of yanking two unrelated economic series out of a hat, forcing them to slow-dance, and then pretending the resulting p-value means something. As always, the official disclaimer: these pairings are chosen by a random number generator with no respect for causality, theory, or the dignity of econometrics. Any relationship you see is almost certainly a cosmic coincidence. Please do not reallocate your retirement fund based on what follows.

We opened the week on Monday with the DISCONTINUED Balance on Services versus the Belize exchange rate, which is already the kind of pairing that makes you want to lie down. The raw regression coughed up a barely-there significant result (R-squared of a towering 9 percent, truly the stuff of legend), and then -- plot twist -- the detrended version actually got slightly stronger and flipped sign. I'd love to tell you this means America's services balance secretly drives the Belizean dollar, but I'd be lying, and so would the data. Tuesday was the week's genuine showstopper: Southern gasoline prices versus 30-year FHA mortgage rates. The raw fit was fine, but the detrended regression held up with an R-squared of 0.63 and a p-value so small it needs scientific notation (1.2 times ten to the minus fifty-one, in case you were bored). Two series that both have something to do with 1980s-90s cost-of-living pressures staying linked even after stripping their trends? That's the rare one that's actually, annoyingly, kind of plausible. Call the Nobel committee. Or don't.

Wednesday brought us corporate machinery cash deposits versus government-sponsored enterprise debt, the textbook trend mirage. The raw regression posted a gaudy 80 percent R-squared and everyone in the room gasped -- and then the detrended version wandered in, posted an R-squared of 2 percent and a p-value of 0.13, and the whole thing evaporated like a puddle in August. Classic. Two things that both go up over decades will always look madly in love until you make them explain themselves without the trend. Thursday's Swedish government consumption prices against Kansas City energy services told the exact same sad story: significant raw, utterly limp (R-squared 0.006, p-value 0.67) once detrended. Friday's corporate bond index versus Middle Atlantic unemployment also folded instantly -- raw hit, detrended nothingburger. Three mirages in a row, which is honestly the project running exactly as designed.

The weekend salvaged some dignity. Saturday paired Indiana property tax collections with Pacific-region nonperforming loans, and the inverse relationship actually survived detrending (R-squared around 0.28, p-value in the billionths), which, if you squint, you could rationalize via the economy-does-well-so-taxes-rise-and-bad-loans-fall story. Squinting optional. Sunday closed with the Singapore dollar exchange rate versus the Atlanta house price index -- raw regression negative, detrended regression flipped positive and stayed significant, which mostly tells you the relationship is unstable enough that you shouldn't build a thesis on it. So this week: two surprisingly sturdy results, a pile of trend-driven illusions, and the enduring reminder that if you torture enough random series pairs, a few will confess to anything. See you next week for more statistically significant nonsense.
